\chapter{Projected subgradient descent}
Recall that we are interested in solving the problem
\begin{equation}
    \min_{x\in C} f(x).
\end{equation}
where $f:C\rightarrow\R$ is convex.
Recall also, that we may always write a constraint problem as above as an unconstraint one via\footnote{potentially extending $f$ arbitrarily outside of $\R^d$}
\begin{equation}
    \min_{x\in \R^d} f(x) + \delta_C(x).
\end{equation}
Let us recall the most basic numerical optimization method for such problems in the case $C=\R^d$: gradient (or \emph{steepest}) descent. We choose an initial value $x_0$ and then perform for $k=0,1,\dots$ the updates
\begin{equation}
    x_{k+1} = x_k-\tau_k \nabla f(x_k)
\end{equation}
where $\tau_k>0$ denotes the step size. As a first \emph{new} algorithm, we consider the subgradient version of steepest descent, for $k=0,1,\dots$
\begin{equation}
    \begin{cases}
        g_k &\in \partial f(x_k)\\
        x_{k+1} &= x_k-\tau_k g_k.
    \end{cases}
\end{equation}
Whenever $C\neq\R^d$ we may quickly run into a problem: It might happen that $x_{k+1}\notin C$ in which case the method is not well defined as the subgradient of $\delta_C(x)$ is empty for $x\notin C$. Therefore, we somehow need to make sure, that each new iterate stays in $C$ for which we will use \emph{projections}.

\begin{theorem}[Projection theorem]
    Let $V$ be an inner product space and $C\subset V$ closed, convex, and non-empty. Then the following function is well-defined
    \begin{equation}
        \begin{aligned}
            \proj_C:\R^d&\rightarrow C\\
            x&\mapsto\arg\min_{y\in C}\|x-y\|.
        \end{aligned}
    \end{equation}
\end{theorem}
\begin{proof}
    We need to show that the minimum is uniquely attained.
    Let $(y_n)_n$ be a minimizing sequence, \ie,
    \begin{equation}
        \lim_n \|y_n-x\| = \inf_{y\in C}\|y-x\|.
    \end{equation}
    Then $(y_n)_n$ is of course bounded. Therefore, there exists a convergent subsequence $y_{n(k)}\rightarrow \hat y$. By closedness of $C$ it follows $\hat y \in C$. Moreover, by continuity of the norm we have
    \begin{equation}
        \|\hat y-x\| = \lim_n \|y_n-x\| = \inf_{y\in C}\|y-x\|.
    \end{equation}
    This yields existence of the minimizer. Regarding uniqueness, let us assume that $y_1\neq y_2$ were both minimizers. Then, using Young's inequality, we find 
    \begin{equation}
        \begin{aligned}
            \| 1/2 (y_1+y_2) - x\|^2 
            =& \frac{1}{4}\|y_1-x\|^2 + \frac{1}{4}\|y_2-x\|^2 + \frac{1}{2}\inner{y_1-x}{y_2-x}\\
            \leq& \frac{1}{4}\|y_1-x\|^2 + \frac{1}{4}\|y_2-x\|^2 + \frac{1}{4}\|y_1-x\|^2 + \frac{1}{4}\|y_2-x\|^2\\
            \leq& \frac{1}{2}\|y_1-x\|^2 + \frac{1}{2}\|y_2-x\|^2\\
            =& \min_{y\in C}\|y-x\|^2.
        \end{aligned}
    \end{equation}
    Note that the above inequality is strict, whenever $y_1-x$ and $y_2-x$ are not colinear and pointed in the same direction. 
    The strict inequality would imply a contradiction to optimality of $y_1$ and $y_2$ as in this case $1/2(y_1+y_2)\in C$ would yield an even smaller value. So it follows that there exists $\lambda\geq 0$ such that
    \[
        y_2-x = \lambda (y_1-x)
    \]
    But since $\|y_1-x\|=\|y_2-x\|$ it follows that $\lambda=1$, \ie, $y_1=y_2$ concluding the proof.
\end{proof}
Equipped with the projection we have a remedy for the problem of \emph{escaping} $C$: the projected subgradient method, for $k=0,1,\dots$
\begin{equation}\label{subgradient:eq:projectedSG}
    \begin{cases}
        g_k &\in \partial f(x_k)\\
        x_{k+1} &= \proj_C( x_k-\tau_k g_k) .
    \end{cases}
\end{equation}
In order to prove convergence of the method we require some auxiliary results.

\begin{lemma}
    Let $C$ be convex. Then it holds
    \begin{equation}
        \inner{x-\proj_C(x)}{y-\proj_C(x)}\leq 0,\quad y\in C.
    \end{equation}
    In particular, if $C$ is a subspace\footnote{this means $C$ fulfills the conditions of a vector space, \cf~\cref{preliminaries:defn:vectorspace}}, it follows
    \begin{equation}
        \inner{x-\proj_C(x)}{y}= 0,\quad y\in C.
    \end{equation}
\end{lemma}
\begin{proof}
    We have for any $y\in C$
    \begin{equation}
        \begin{aligned}
            \|\proj_C(x)-x\|^2\leq& \|y-x\|^2\\
            =& \|y - \proj_C(x) + \proj_C(x)-x\|^2\\
            =& \|y-\proj_C(x)\|^2 + 2 \inner{y-\proj_C(x)}{\proj_C(x)-x} + \|\proj_C(x)-x\|^2.
        \end{aligned}
    \end{equation}
    As a consequence we have
    \begin{equation}
        \begin{aligned}
            0\leq \|y-\proj_C(x)\|^2 + 2 \inner{y-\proj_C(x)}{\proj_C(x)-x}
        \end{aligned}
    \end{equation}
    Now let $z\in C$ arbitrary and plug in $y=\proj_C(x) + t(z-\proj_C(x))$ for $t\in (0,1)$ above. Then we find
    \begin{equation}
        \begin{aligned}
            0\leq t^2 \|z-\proj_C(x)\|^2 + 2 t\inner{z-\proj_C(x)}{\proj_C(x)-x}.
        \end{aligned}
    \end{equation}
    Dividing by $t>0$ we find
    \begin{equation}
        \begin{aligned}
            0\leq t \|z-\proj_C(x)\|^2 + 2 \inner{z-\proj_C(x)}{\proj_C(x)-x}.
        \end{aligned}
    \end{equation}
    Letting $t\rightarrow 0$ yields the desired result. The equality in the case of a subspace is left as an exercise.
\end{proof}


\begin{prop}
    The projection onto a convex set is nonexpansive, \ie, Lipschitz continuous with Lipschitz constant one, that is
    \begin{equation}
        \|\proj_C(x_1)-\proj_C(x_2)\|\leq \|x_1-x_2\|.
    \end{equation}
\end{prop}
\begin{proof}
    Assume without loss of generality $\|\proj_C(x_2)-\proj_C(x_1)\|^2 \neq 0$ as otherwise there is nothing to prove. We find
    \begin{equation}
        \begin{aligned}
            \|\proj_C(x_2)-\proj_C(x_1)\|^2 
            =& \inner{\proj_C(x_2)-\proj_C(x_1)}{\proj_C(x_2)-x_1+x_1-\proj_C(x_1)}\\
            =& \inner{\proj_C(x_2)-\proj_C(x_1)}{\proj_C(x_2)-x_1} \\
            &+ \inner{\proj_C(x_2)-\proj_C(x_1)}{x_1-\proj_C(x_1)}\\
            \leq& \inner{\proj_C(x_2)-\proj_C(x_1)}{\proj_C(x_2)-x_1}\\
            =& \inner{\proj_C(x_2)-\proj_C(x_1)}{\proj_C(x_2)-x_2+x_2-x_1}\\
            \leq& \inner{\proj_C(x_1)-\proj_C(x_2)}{x_2-\proj_C(x_2)}\\
            &+ \inner{\proj_C(x_2)-\proj_C(x_1)}{x_2-x_1}\\
            \leq& \inner{\proj_C(x_2)-\proj_C(x_1)}{x_2-x_1}\\
            \leq& \|\proj_C(x_2)-\proj_C(x_1)\|\|x_2-x_1\|
        \end{aligned}
    \end{equation}
    Dividing by $\|\proj_C(x_2)-\proj_C(x_1)\|\neq 0$ yields the result.
\end{proof}


\begin{lemma}[Fundamental inequality of the projected subgradient method]\label{lemma:subgradient:fundamental}
    Let $x_k$ denote the iterates of the projected subgradient method~\eqref{subgradient:eq:projectedSG} and $g_k\in\partial f(x_k)$. Then it holds true that 
    \begin{equation}
        \|x_{k+1}-x^*\|^2 
            \leq \|x_k- x^*\|^2 - 2\tau_k (f(x_k)-f(x^*)) + \tau_k^2\|g_k\|^2
    \end{equation}
\end{lemma}
\begin{proof}
    By definition of the subgradient we have $f(x_k)+ \inner{g_k}{x^*-x_k}\leq f(x^*)$. It follows
    \begin{equation}\label{eq:subgradient:fundamental}
        \begin{aligned}
            \|x_{k+1}-x^*\|^2 
            =& \|\proj_C(x_k-\tau_kg_k) - \proj_C(x^*)\|^2\\
            \leq& \|x_k-\tau_k g_k - x^*\|^2\\
            \leq& \|x_k- x^*\|^2 - 2\tau_k \inner{g_k}{x_k-x^*} + \tau_k^2\|g_k\|^2\\
            \leq& \|x_k- x^*\|^2 - 2\tau_k (f(x_k)-f(x^*)) + \tau_k^2\|g_k\|^2
        \end{aligned}
    \end{equation}
\end{proof}
A natural thing to do is now to use the previous result to deduce an optimal step size. More precisely, minimizing the right-hand side of~\eqref{eq:subgradient:fundamental} with respect to $\tau_k$ we find that the optimal step-size is
\begin{equation}
    \tau_k = \frac{f(x_k)-f(x^*)}{\|g_k\|^2}.
\end{equation}
whenever $g_k\neq 0$, and arbitrary (\eg, $\tau_k=1$) whenever $g_k=0$ as in the latter case the iteration stagnates anyway.
This step size is referred to as Polyak's step size rule. We obtain the following result.

\begin{theorem}
    Assume that $f:C\rightarrow \R$ is Lipschitz continuous with Lipschitz constant $L$. With Polyak's step size rule 
    \begin{equation}
    \tau_k = \frac{f(x_k)-f(x^*)}{\|g_k\|^2}
\end{equation}
    the projected subgradient method~\eqref{subgradient:eq:projectedSG} satisfies
    \begin{enumerate}
        \item $\|x_{k+1}-x^*\|^2\leq \|x_{k}-x^*\|$ with strict inequality whenever $g_k\neq 0$,
        \item $f(x_K)\rightarrow f(x^*)$ as $k\rightarrow\infty$, and
        \item $f_{\mathrm{best}}^N-f(x^*)\leq \frac{L\|x_1-x^*\|}{\sqrt{N}}$ where $f_{\mathrm{best}}^N = \min_{k=1,\dots,N}f(x_k)$.
    \end{enumerate}
\end{theorem}
\begin{proof}
    Without loss of generality, we may assume $f(x^*)=0$. Otherwise simply consider $\tilde{f}(x) \coloneq f(x)-f(x^*)$. By~\cref{lemma:subgradient:fundamental} and the step size rule we obtain
    \begin{equation}
        \begin{aligned}
            \|x_{k+1}-x^*\|^2 
            \leq& \|x_{k} - x^*\|^2 - 2\tau_k f(x_k) +\tau_k^2\|g_k\|^2\\
            =& \|x_{k} - x^*\|^2 - a_k\\
        \end{aligned}
    \end{equation}
    where $a_k=\frac{f(x_k)^2}{\|g_k\|^2}$ if $g_k\neq 0$ and zero else.
    By rearranging and summing over $k$ we find
    \begin{equation}
        \sum_{k=1}^{N-1} a_k \leq \|x_1-x^*\|^2 - \|x_N-x^*\|^2\leq \|x_1-x^*\|^2
    \end{equation}
    Lipschitz continuity of $f$ implies that $\|g_k\|\leq L$ (exercise!) so that we obtain
    \begin{equation}
        \sum_{k=1}^{N-1} f(x_k)^2\leq \sum_{k=1}^{N-1} L^2 a_k \leq L^2 \|x_1-x^*\|^2
    \end{equation}
    Since the above is true for any $N$ we may let $N\rightarrow \infty$. As a consequence we find $f(x_k)\rightarrow 0$. Moreover,
    \begin{equation}
        N (f_{\mathrm{best}}^N)^2\leq \sum_{k=1}^{N} f(x_k) \leq L^2\|x_1-x^*\|^2
    \end{equation}
    \ie, 
    \begin{equation}
        f_{\mathrm{best}}^N\leq \frac{L\|x_1-x^*\|}{\sqrt{N}}.
    \end{equation}
\end{proof}

\begin{cor}
    We require $N=\Oc(\frac{L^2\|x_1-x^*\|^2}{\epsilon^2})$ iterations in order to reach accuracy $f_{\mathrm{best}}^N - f(x^*)\leq \epsilon$.
\end{cor}

\begin{rem}
    Recall that standard smooth gradient methods achieve complexity $\Oc(\frac{1}{\epsilon})$.
\end{rem}

While the convergent result using Polyak's step size rule is theoretically interesting and yields a best-case complexity result, in practice we will not be able to compute the step size as it requires knowledge of $f(x^*)$. In the following, we provide a more general convergence result.

\begin{theorem}
    Let $f$ be $L$-Lipschitz and assume the step sizes satsify
    \[
        \frac{\sum_{k=1}^n\tau_k}{\sum_{k=1}^n\tau_k^2}\rightarrow \infty
    \]
\end{theorem}
as $n\rightarrow\infty$. Then it holds $f^n_{\mathrm{best}}\rightarrow f(x^*)$.
\begin{proof}
    We assume again for simplicity $f(x^*)=0$.
    Using again $\|g_k\|\leq L$ and summing the fundamental inequality~\cref{lemma:subgradient:fundamental} over $k=1,\dots, N$ yields
    \begin{equation}
        2\sum_{k=1}^N\tau_k f(x_k) \leq \|x_1-x^*\|^2 + L^2 \sum_{k=1}^N \tau_k^2 
    \end{equation}
    Let us denote
    \[
        \sigma_n = \frac{\sum_{k=1}^n\tau_k}{\sum_{k=1}^n\tau_k^2}
    \]
    It follows
    \begin{equation}
        \sigma_n f^n_{\mathrm{best}} 
        \leq \frac{\sum_{k=1}^n\tau_kf^n_{\mathrm{best}}}{\sum_{k=1}^n\tau_k^2}
        \leq \frac{\sum_{k=1}^n\tau_kf(x_k)}{\sum_{k=1}^n\tau_k^2}
        \leq \frac{\|x_1-x^*\|^2}{2\sum_{k=1}^n\tau_k^2} + L^2/2.
    \end{equation}
    Since the right-hand side is bounded as $n\rightarrow\infty$, $\sigma_n\rightarrow\infty$ implies $f^n_{\mathrm{best}}\rightarrow 0$.
\end{proof}

Variants of the convergence results can moreover be obtained if we assume, \eg, that $C$ is compact and we refer to~\cite{beck2017first} for details.

The convergence can be improved and also \emph{transferred} to the iterates $(x_k)_k$ is we additionally assume that $f$ is strongly convex.


\begin{theorem}
    Assume that $f$ is $\mu$ strongly convex and $L$-Lipschitz over $C$\footnote{This implies that $C$ is bounded as every strongly convex function on an unbounded domain admits an unbeunded gradient!}. Then with the step size choice $\tau_k = \frac{2}{\mu k}$ the projected subgradient method~\eqref{subgradient:eq:projectedSG} satisfies
    \begin{enumerate}
        \item $f^n_{\mathrm{best}}-f(x^*)\leq \frac{L^2}{\mu (n-1)}$ and
        \item $\|x_n-x^*\|\leq \frac{2L}{\mu \sqrt{n-1}}$.
    \end{enumerate}
\end{theorem}
\begin{proof}
    As always, without loss of generality $f(x^*)=0$. Under strong convexity the fundamental inequality can in fact be improved to 
    \begin{equation}
        \begin{aligned}
            \|x_{k+1}-x^*\|^2 
            &\leq (1-\mu\tau_k)\|x_k- x^*\|^2 - 2\tau_k f(x_k) + \tau_k^2\|g_k\|^2\\
            &\leq (1-\mu\tau_k)\|x_k- x^*\|^2 - 2\tau_k f(x_k) + \tau_k^2L^2
        \end{aligned}
    \end{equation}
    Rearranging and dividing by $2\tau_k$ yields
    \begin{equation}
        \begin{aligned}
            f(x_k)
            \leq -\frac{\tau_k^{-1}}{2}\|x_{k+1}-x^*\|^2 + \frac{\tau_k^{-1}-\mu}{2}\|x_k- x^*\|^2  + \frac{\tau_k}{2}L^2.
        \end{aligned}
    \end{equation}
    Inserting the step size choice $\tau_k = \frac{2}{\mu k}$ then yields
    \begin{equation}
        \begin{aligned}
            f(x_k)
            \leq -\frac{\mu k}{4} \|x_{k+1}-x^*\|^2 + \frac{\mu(k-2)}{4}\|x_k- x^*\|^2  + \frac{1}{\mu k}L^2.
        \end{aligned}
    \end{equation}
    Multyplying by $k-1$ and summing over $k$ yields
    \begin{equation}
        \sum_{k=1}^{n} (k-1)f(x_k) \leq -\frac{\mu n(n-1)}{4} \|x_{n+1}-x^*\|^2 + \frac{1}{\mu}\sum_{k=1}^{n}\frac{k-1}{k}L^2.
    \end{equation}
    Using that $\sum_{k=1}^{n} (k-1) = n(n-1)/2$ and that $\sum_{k=1}^{n}\frac{k-1}{k}\leq n$ we obtain
    \begin{equation}\label{subgradient:eq:strongly}
        n(n-1)/2 f^n_{\mathrm{best}}\leq \sum_{k=1}^{n} (k-1)f(x_k) = -\frac{\mu n(n-1)}{4} \|x_{n+1}-x^*\|^2 + \frac{1}{\mu} n L^2.
    \end{equation}
    Thus
    \begin{equation}
        f^n_{\mathrm{best}} \leq \frac{"L^2}{\mu(n-1)}.
    \end{equation}
    On the other hand-side, rearranging \eqref{subgradient:eq:strongly} and using that $f^n_{\mathrm{best}}\geq 0$ yields
    \begin{equation}
        \frac{\mu n(n-1)}{4} \|x_{n+1}-x^*\|^2 \leq \frac{1}{\mu} n L^2,
    \end{equation}
    hence, $\|x_{n+1}-x^*\|^2\leq \frac{4L^2}{\mu^2 (n-1)}$.
\end{proof}

